\expandafter\def\csname HistoricalNoteText27\endcsname{Dieses Werk über die Bestimmung der Planetenbahnen wird nun bald erscheinen. Dr. \textsc{Gauss} schreibt uns, dass die \textsc{Perthes}'sche Buchhandlung in Hamburg den Verlag davon übernommen habe. Das Manuscript ist zwar ganz vollendet; da Dr. \textsc{Gauss} aber dem Wunsche des Herr Verlegers zu Folge, um das Werk auch solchen Ausländern lesbar zu machen, die der deutschen Sprache nicht kundig sind, dasselbe nun lateinisch herausgibt, so verursacht dies noch einigen Aufenthalt; in ein paar Monaten wird aber der Druck anfangen können.}
\clearpage
\fancyhead[L]{\small GAUSS AN BODE.}
\hypertarget{native-289-lower}{}
\begin{flushright}
Braunschweig, den 23. August 1807.
\end{flushright}

\par\indent Verzeihen Ew. es gütigst, dass ich Ihnen seit der Mittheilung meiner ersten Elemente der Vesta
noch keine weitern Nachrichten von meinen spätern Untersuchungen über diesen Gegenstand eingesandt habe.
\SourceLeaf{290}%
\par\indent Eine in diesem Sommer nach Bremen gemachte Reise, und eine mir nachher zugestossene Krankheit haben mich, sowie in meinen Arbeiten überhaupt, so auch in meiner Correspondenz etwas zurückgesetzt: jetzt eile ich aber um so mehr das Versäumte nachzuholen, damit Sie noch nach Gefallen von meinen Mittheilungen in Ihrem neuen Jahrbuche Gebrauch machen können.

\par\indent Meine zweiten Elemente der \VestaBodyTwoNinety{} und deren Uebereinstimmung mit sehr vielen Beobachtungen werden Ihnen bereits aus der M.\,C. bekannt sein.

\par\indent Zur Berechnung der \emph{dritten} Elemente habe ich ausser den frühern Beobachtungen, noch die Pariser und Mailänder Meridianbeobachtungen und die letzten Bremer und Lilienthaler --- welche bis zum 11. Juli reichen --- benutzt. Hier diese dritten Elemente:

\begin{center}\begingroup\fontsize{9}{11}\selectfont
\begin{tabular}{ll}
\shortstack[l]{Epoche der Länge 1807. März 31.\\$0^{\text{u}}$ in Bremen} .\ .\ . & $192^\circ\,23'\,30''1$ \\
Sonnenferne .\ .\ .\ .\ .\ .\ .\ .\ .\ .\ .\ .\ .\ .\ .\ . & $69\;50\;31.9$ \\
Knoten .\ .\ .\ .\ .\ .\ .\ .\ .\ .\ .\ .\ .\ .\ .\ .\ .\ . & $103\;18\;28.0$
\end{tabular}\endgroup\end{center}

\noindent Beide für obige Epoche und siderisch ruhend vorausgesetzt.

\begin{center}\begingroup\fontsize{9}{11}\selectfont
\begin{tabular}{ll}
Neigung der Bahn .\ .\ .\ .\ .\ .\ .\ .\ .\ .\ .\ .\ . & $7\;8\;10.7$ \\
Excentricität .\ .\ .\ .\ .\ .\ .\ .\ .\ .\ .\ .\ .\ .\ . & $0.0855050$ \\
Logarithmus der halben Axe .\ .\ .\ .\ .\ .\ .\ .\ .\ . & $0.3720160$ \\
Tägliche mittlere tropische Bewegung .\ .\ .\ .\ .\ . & $981''7087$
\end{tabular}\endgroup\end{center}

\par\indent Ich füge hier zugleich die nach diesen Elementen berechnete Ephemeride für das nächste Jahr bei.

[Vom 2. Mai 1808 bis 18. März 1809 mit viertägigen Zwischenräumen für Mitternacht in Seeberg sind AR. und Decl. auf Minuten genau gegeben.]

\clearpage\fancyhead[L]{\small VESTA.}
\begin{center}
VESTA.
\end{center}

\begin{center}
\rule{0.6\linewidth}{0.4pt}\\
Monatliche Correspondenz. Band XV. S. 590–600. 1807 Juni.\\
\rule{0.6\linewidth}{0.4pt}
\end{center}

--- --- ---

\par\indent Herr Dr. \textsc{Gauss} in Braunschweig machte folgende zwei Beobachtungen, die sich auf zahlreiche Vergleichungen mit dem \textsc{Piazzi}'schen Stern, AR. $180^\circ\,51'\,27''5$, nördl. Abweich. $12^\circ\,35'\,14''0$ gründen.

\begin{center}\begingroup
\fontsize{10}{12}\selectfont
\setlength{\tabcolsep}{2.5pt}\renewcommand{\arraystretch}{1.15}
\begin{tabular}{c|c|c|c}
1807. & \shortstack{Mittl. Zeit in\\Braunschweig} & \shortstack{Scheinb. gerade\\Aufsteigung \VestaRATwoNinety{}} & \shortstack{Scheinb. nördl.\\Abweichung \VestaDecTwoNinety{}} \\ \hline
April 6 & $8^{\text{u}}\,22^{\text{m}}\,45^{\text{s}}$ & $182^\circ\,21'\,54''6$ & $12^\circ\,28'\,0''9$ \\
8 & $8\ \ 28\ \ 49$ & $181\ \ 56\ \ 51.0$ & $12\ \ 35\ \ 45.8$
\end{tabular}\endgroup\end{center}

\par\indent Es war zu erwarten, dass der bewunderungswürdige Dr. \textsc{Gauss}, welcher sich so ausschliesslich Verdienste um alle neu entdeckte kleine Planeten erworben hat, diese auch bei der \emph{Vesta} einzig behaupten würde. Allein diesmal hat sich dieser grosse Geometer selbst übertroffen. „Unsern vortrefflichen \textsc{Gauss} habe ich schon oft bewundert, (schreibt uns Dr. \textsc{Olbers} unterm 20. April) aber diesmal hat er mich doch in Erstaunen gesetzt. Er erhielt meine Beobachtung der \emph{Vesta} vom 17. April am 20. April Abends um $9\tfrac{1}{2}$ Uhr und am 21. April Nachmittags um 5 Uhr konnte er mir schon die berechneten elliptischen Elemente und die Vergleichung mit allen Beobachtungen schicken. Zu diesem allen-
\SourceLeaf{291}%
hat er nur zehn Stunden gebraucht. Zugleich ein Beweis der grossen Geschmeidigkeit seiner neuen so sehr vervollkommneten Methode\HistoricalNote{27}, Planetenbahnen zu bestimmen.“

\par\indent Dr. \textsc{Gauss} schreibt uns hierüber selbst folgendes: „Die erste Bestimmung der Bahn unternahm ich am 20. April, als ich Dr. \textsc{Olbers} Beobachtung vom 17. April erhalten hatte, wodurch die Zwischenzeit auf 19 Tage angewachsen war. Die Resultate übersandte ich am 21. April an Dr. \textsc{Olbers}. Da aber die Zwischenzeit noch so klein, die Beobachtungen fast alle nur mit dem Kreismicrometer gemacht waren und ausserdem die kleine Neigung der Bahn dieses Planeten auch noch dazu beiträgt, die Resultate verhältnissmässig weniger zuverlässig zu machen, so wagte ich es noch nicht, sie schon für eine wirkliche Annäherung zur Wahrheit auszugeben und hielt sie daher mit dem Vorsatz zurück, sie erst noch einmal durch spätere Beobachtungen zu verbessern. Erst gestern (27. April) sind mir neuere Beobachtungen zugekommen, womit ich nun gleich eine neue Berechnung zu machen dachte; allein zu meinem grossen Vergnügen finde ich, dass meine ersten Resultate damit noch so genau übereinstimmen, dass ich eigentlich mit Zuverlässigkeit noch gar nichts daran zu ändern wüsste. Indem ich also diese neue Rechnung so lange aufschiebe, bis der Fehler zu einer entschiedenen merklichen Grösse angewachsen sein wird, warte ich nun nicht länger, Ihnen die I. Elemente der \emph{Vesta} mitzutheilen.

\begin{center}\begingroup\fontsize{9}{11}\selectfont
\begin{tabular}{ll}
\shortstack[l]{Epoche der Länge 1807 März 29.\\$12^{\text{u}}$ M.\,Z. in Bremen} . & $193^\circ\,8'\,4''6$ \\
Sonnenferne .\ .\ .\ .\ .\ .\ .\ .\ .\ .\ .\ .\ .\ .\ .\ . & $69^\circ\;7\;40.9$ \\
Aufsteigender Knoten .\ .\ .\ .\ .\ .\ .\ .\ .\ .\ .\ .\ . & $103\;8\;36.2$
\end{tabular}\endgroup\end{center}

\noindent Alle drei vom mittlern Aequinoctium gezählt und die beiden letztern als tropisch ruhend betrachtet.

\begin{center}\begingroup\fontsize{9}{11}\selectfont
\begin{tabular}{ll}
Tägliche mittlere Bewegung .\ .\ .\ .\ .\ .\ .\ .\ .\ .\ .\ . & $978''909$ \\
Excentricität .\ .\ .\ .\ .\ .\ .\ .\ .\ .\ .\ .\ .\ .\ .\ .\ .\ . & $0.097505$ \\
Logarithm der halben grossen Axe .\ .\ .\ .\ .\ .\ .\ .\ . & $0.3728428$ \\
Neigung der Bahn .\ .\ .\ .\ .\ .\ .\ .\ .\ .\ .\ .\ .\ .\ . & $7^\circ\;5'\;49''5$
\end{tabular}\endgroup\end{center}

\par\indent „Hier nun die Vergleichung mit sämmtlichen bekannt gewordenen Beobachtungen.

\begin{center}\begingroup
\fontsize{8}{10}\selectfont
\setlength{\tabcolsep}{2.5pt}\renewcommand{\arraystretch}{1.15}
\begin{tabular}{c|c|c|c|c|c|c}
1807 & \shortstack{Mittlere Zeit\\am\\Beob. Orte} & \shortstack{Berechn. gerade\\Aufsteigung\\\VestaRATwoNinetyOne{}} & \shortstack{Unter-\\schied.} & \shortstack{Berechn. nördl.\\Abweichung\\\VestaDecTwoNinetyOne{}} & \shortstack{Unter-\\schied.} & Beobachter. \\ \hline
März 29 & $8^{\text{u}}\,21^{\text{m}}\,20^{\text{s}}$ & $184^\circ\,8'\,54''4$ & $-1''4$ & $.\ .\ .\ .\ .$ & $.\ .\ .$ & Olbers \\
        & $10\ \ 31\ \ 16$ & $184\ \ 7\ \ 40.2$ & $-10.6$ & $11^\circ\,47'\,52''4$ & $+5''4$ & \\
30 & $8\ \ 34\ \ 53$ & $183\ \ 55\ \ \ 3.0$ & $+6.2$ & $.\ .\ .\ .\ .$ & $.\ .\ .$ & \\
   & $8\ \ 44\ \ \ 8$ & $183\ \ 54\ \ 57.7$ & $+11.9$ & $11\ \ 53\ \ 18.6$ & $+7.6$ & \\
   & $12\ \ 33\ \ 17$ & $183\ \ 52\ \ 44.2$ & $+3.4$ & $11\ \ 54\ \ 14.7$ & $-12.3$ & \\
April 1 & $9\ \ 50\ \ \ 0$ & $183\ \ 27\ \ \ 3.7$ & $+0.9$ & $12\ \ \ 4\ \ 40.7$ & $-11.3$ & \\
   & $12\ \ \ 4\ \ 35$ & $183\ \ 25\ \ 47.8$ & $+23.6$ & $12\ \ \ 5\ \ \ 9.6$ & $-22.1$ & Bessel \\
2 & $8\ \ 21\ \ \ 1$ & $183\ \ 14\ \ 23.8$ & $+6.0$ & $12\ \ \ 9\ \ 34.7$ & $-12.3$ & Olbers \\
3 & $8\ \ 16\ \ 49$ & $183\ \ \ 1\ \ \ 3.6$ & $+2.8$ & $12\ \ 14\ \ 33.6$ & $-6.4$ & \\
4 & $9\ \ \ 9\ \ 43$ & $182\ \ 47\ \ 20.9$ & $+0.1$ & $12\ \ 19\ \ 29.1$ & $-15.9$ & \\
5 & $11\ \ 17\ \ \ 3$ & $182\ \ 33\ \ 10.7$ & $-0.2$ & $12\ \ 24\ \ 21.6$ & $+2.5$ & Harding \\
6 & $8\ \ 22\ \ 45$ & $182\ \ 21\ \ 48.5$ & $-6.1$ & $12\ \ 28\ \ \ 6.6$ & $+5.7$ & Gauss \\
  & $11\ \ 12\ \ 15$ & $182\ \ 20\ \ 21.6$ & $-12.2$ & $12\ \ 28\ \ 34.6$ & $+45.2$ & Harding \\
8 & $8\ \ 21\ \ 37$ & $181\ \ 56\ \ 42.9$ & $-0.1$ & $12\ \ 35\ \ 50.1$ & $-8.0$ & Olbers \\
  & $8\ \ 28\ \ 49$ & $181\ \ 56\ \ 42.2$ & $-8.8$ & $12\ \ 35\ \ 50.1$ & $+4.3$ & Gauss \\
  & $11\ \ \ 2\ \ 43$ & $181\ \ 55\ \ 22.5$ & $+3.3$ & $12\ \ 36\ \ 13.6$ & $+1.4$ & Harding
\end{tabular}\endgroup\end{center}
\begin{center}\begingroup
\fontsize{8}{10}\selectfont
\setlength{\tabcolsep}{2.5pt}\renewcommand{\arraystretch}{1.15}
\begin{tabular}{c|c|c|c|c|c|c}
\SourceLeaf{292}%
1807 & \shortstack{Mittlere Zeit\\am\\Beob.-Orte} & \shortstack{Berechn. gerade\\Aufsteigung\\\VestaRATwoNinetyTwo{}} & \shortstack{Unter-\\schied.} & \shortstack{Berechn. nördl.\\Abweichung\\\VestaDecTwoNinetyTwo{}} & \shortstack{Unter-\\schied.} & Beobachter. \\ \hline
April 9 & $8^{\text{u}}\,30^{\text{m}}\,46^{\text{s}}$ & $181^\circ\,44'\,25''6$ & $+8''6$ & $12^\circ\,39'\,19''8$ & $+5''7$ & Olbers \\
   & $10\ \ 57\ \ 58$ & $181\ \ 43\ \ 13.8$ & $+11.5$ & $12\ \ 39\ \ 39.6$ & $+3.5$ & Harding \\
12 & $8\ \ 27\ \ 20$ & $181\ \ \ 9\ \ 25.6$ & $-8.4$ & $12\ \ 48\ \ \ 7.4$ & $-8.7$ & Olbers \\
13 & $10\ \ 39\ \ 14$ & $180\ \ 57\ \ 28.3$ & $+11.3$ & $12\ \ 50\ \ 41.5$ & $-1.5$ & Bode \\
14 & $8\ \ 19\ \ 51$ & $180\ \ 47\ \ 36.3$ & $-21.7$ & $12\ \ 52\ \ 37.2$ & $-10.9$ & Olbers \\
   & $10\ \ 34\ \ 33$ & $180\ \ 46\ \ 39.2$ & $-0.7$ & $12\ \ 52\ \ 44.4$ & $+4.7$ & Harding \\
17 & $8\ \ 12\ \ 35$ & $180\ \ 17\ \ 21.6$ & $-5.2$ & $12\ \ 57\ \ 18.3$ & $+2.3$ & Olbers \\
23 & $8\ \ 30\ \ 52$ & $179\ \ 26\ \ 39.7$ & $-1.3$ & $12\ \ 59\ \ 20.4$ & $+0.4$ & \\
25 & $9\ \ 53\ \ 49$ & $179\ \ 13\ \ \ 0.6$ & $-18.4$ & $12\ \ 57\ \ 52.4$ & $-5.6$ & \\
   & $9\ \ 49\ \ 26$ & $179\ \ 12\ \ 57.1$ & $+8.1$ & $12\ \ 57\ \ 51.9$ & $-14.1$ &
\end{tabular}\endgroup\end{center}

\par\indent Diese Übereinstimmung ist so gut, als sie sich bei diesen Beobachtungen nur erwarten lässt, und da nunmehr die Zwischenzeit 27 Tage beträgt, so ist Dr. \textsc{Gauss} der Meinung, dass diese Elemente wenigstens eine Idee von den Haupt-Dimensionen der Bahn geben. Die Umlaufszeit der Bahn wäre demnach bei der \emph{Vesta} viel schneller, als bei den andern kleinen Planeten, die Neigung der Bahn von allen am kleinsten, die Excentricität mässig; der letzte Umstand und die grosse Helligkeit des Planeten geben Dr. \textsc{Gauss} sehr grosse Hoffnung, dass wir von diesem neuen Planeten ältere Beobachtungen auffinden werden, vielleicht selbst bei \textsc{Flamsteed}. Sobald die Elemente nur mit etwas grösserer Zuverlässigkeit bekannt sein werden, wird Dr. \textsc{Gauss} darüber Nachsuchungen anstellen.

\par\indent Unterm 18. Mai schreibt uns dieser unvergleichliche Astronom: „Meine erste Ihnen schon vor einiger Zeit übersandten Elemente der \emph{Vesta}, die auf 19tägige Beobachtungen gegründet waren, stimmten bis gegen das Ende des Aprils fortwährend sehr gut mit den Beobachtungen überein: im Mai aber fing der Unterschied an, einigermaassen merklich zu werden und war am 6. Mai bis zu $+33''$ in gerader Aufsteigung und $-41''$ in Declination angewachsen. Nicht gerade, weil eine Verbesserung der Elemente nun schon nöthig gewesen wäre, sondern hauptsächlich, weil ich einiger mir verdächtig gewordener Sterne aus älteren Catalogen wegen ungeduldig war, eine Idee davon zu erhalten, wie grosse Änderungen die Elemente etwa noch erfordern würden, unternahm ich am 9. Mai, als ich Dr. \textsc{Olbers} Beobachtung vom $6^{\text{ten}}$ erhalten hatte, eine zweite Berechnung der Bahn, von der ich hier die Resultate mitzutheilen die Ehre habe.

\Needspace{16\baselineskip}
\begin{center}
II. \emph{Elemente der Vesta}
\end{center}

\begin{center}
auf 38tägige Beobachtungen gegründet.
\end{center}

\begin{center}\begingroup\fontsize{9}{11}\selectfont
\begin{tabular}{ll}
\shortstack[l]{Epoche der Länge, 1807, März 29. $12^{\text{u}}$\\M.\,Z. in Bremen} .\ .\ .\ .\ .\ . & $192^\circ\,9'\,53''9$ \\
Sonnenferne .\ .\ .\ .\ .\ .\ .\ .\ . & $69\;57\;52$ \\
Aufsteigender Knoten .\ .\ .\ .\ .\ . & $103\;18\;34$ \\
Tägliche mittlere tropische Bewegung  & $980''707$ \\
Excentricität .\ .\ .\ .\ .\ .\ .\ .\ . & $0.0872230$ \\
Logarithm der mittl. Entfernung . & $0.3723521$ \\
Neigung der Bahn .\ .\ .\ .\ .\ .\ . & $7^\circ\,8'\,6''8$
\end{tabular}\endgroup\end{center}

\par\indent „Diese Elemente stimmen noch bis zum $11^{\text{ten}}$, als so weit meine erhaltenen Beobachtungen reichen, auf's schönste mit diesen überein; ich habe bisher zusammen 68 Beobachtungen verglichen; die Übereinstimmung ist durchgehends so, dass sich nirgends eine entschiedene Differenz zwischen dem berechneten und dem wahren Orte angeben lässt.“

\par\indent „Meine Hoffnung, dass wir Beobachtungen der \emph{Vesta} in der \emph{Historia coelestis} antreffen werden, ist noch immer sehr gross, aber bei denjenigen vermissten Sternen, die ich gleich anfangs als verdächtig angezeichnet hatte, ist die Hoffnung grösstentheils verschwunden. Blos über 91 \emph{Virginis} nach
\SourceLeaf{293}%
\textsc{Flamsteed} will ich noch nicht absprechen, die Länge würde sich mit dem Orte der \emph{Vesta} sehr gut vereinigen lassen, wenn die mittlere tägliche Bewegung nur etwa $6''$ grösser wäre, als nach den II. Elementen; aber die Breite \emph{scheint} etwas zu gross. Bald werden wir hierüber bestimmter urtheilen können.“

--- --- ---
\hfill [\textsc{von Zach}.]
